% Skriptaufgaben -- Analysis (Modul 61211), Lektion 1
% Quelle: die im Lehrtext-Skript (61211-01-S.pdf) eingebetteten UEBUNGSAUFGABEN
%   "Ü 1.x.y" am Ende jedes Abschnitts, samt der offiziellen Loesungen aus dem
%   Loesungsteil "Loesungen der Aufgaben zu 1.x" desselben Skripts.
% Format: \lernkarte[id]{Vorderseite: Ü-Nummer + Aufgabe}{Rueckseite: kurze Loesung}
% id-Schema: 61211-1-sa-NN  ("sa" = Skriptaufgabe)
%
% --- Abschnitt 1.1 Reelle Zahlen -----------------------------------------
\lernabschnitt{1.1 Reelle Zahlen}

\lernkarte[61211-1-sa-01]{%
  Übung 1.1.1\\[4pt]
  Beweisen Sie 1.1.2: Für $a,b\in\mathbb{R}$ (bzw. $a\neq0$) besitzt
  $a+x=b$ bzw. $a\,y=b$ genau eine Lösung.
}{%
  \textbf{Existenz:} $x_0:=-a$ löst $a+x_0=0$, dann $x:=x_0+b=-a+b$.\\
  \textbf{Eindeutigkeit:} aus $a+x=b=a+\tilde x$ durch Addition von $-a$
  (Kürzen) $x=\tilde x$. Für $a\,y=b$ analog mit $a^{-1}$ ($a\neq0$).
}

\lernkarte[61211-1-sa-02]{%
  Übung 1.1.2\\[4pt]
  Berechnen Sie zu $y=\begin{pmatrix}a&b\\-b&a\end{pmatrix}\neq\mathbf{0}$
  die Inverse $y^{-1}$ in $\mathbb{C}$ und zeigen Sie für
  $\mathbf{1}=\begin{pmatrix}1&0\\0&1\end{pmatrix}$,
  $\mathbf{i}=\begin{pmatrix}0&1\\-1&0\end{pmatrix}$:
  $y=a\mathbf{1}+b\mathbf{i}$ und $\mathbf{i}^2=-\mathbf{1}$.
}{%
  \[ y^{-1}=\frac{1}{a^2+b^2}\begin{pmatrix}a&-b\\b&a\end{pmatrix}, \]
  denn $a^2+b^2\neq0$. $\mathbf{i}^2=-\mathbf{1}$ und
  $a\mathbf{1}+b\mathbf{i}=y$ durch Nachrechnen.
}

\lernkarte[61211-1-sa-03]{%
  Übung 1.1.3\\[4pt]
  Sei $(\mathbb{K},v)$ ein linear geordneter Körper. Zeigen Sie
  $0\,v\,a^2$ für jedes $a\in\mathbb{K}$. Können $\mathbb{F}_2$ und
  $\mathbb{C}$ linear geordnet werden?
}{%
  $0\,v\,a^2$ gilt stets: Fallunterscheidung $0\,v\,a$ bzw. $a\,v\,0$,
  dann mit $a$ bzw. $-a$ multiplizieren ($(-a)(-a)=a^2$).\\
  $\mathbb{F}_2$: nein ($1\,v\,0$ und $0\,v\,1$ $\Rightarrow 1=0$).
  $\mathbb{C}$: nein ($\mathbf{i}^2=-1\Rightarrow 0\,v\,-1$, Widerspruch).
}

\lernkarte[61211-1-sa-04]{%
  Übung 1.1.4\\[4pt]
  Sei $M\subseteq\mathbb{R}$ nichtleer. Zeigen Sie: $M$ ist genau dann
  beschränkt, wenn $|M|:=\{\,|x|\mid x\in M\,\}$ nach oben beschränkt ist.
}{%
  \textbf{$\Rightarrow$:} $s\le x\le S \Rightarrow |x|\le\max\{S,-s\}$.\\
  \textbf{$\Leftarrow$:} $|x|\le K \Rightarrow -K\le x\le K$, also
  $M$ beschränkt.
}

\lernkarte[61211-1-sa-05]{%
  Übung 1.1.5\\[4pt]
  Sei $A\subseteq\mathbb{R}$ nichtleer, nach unten beschränkt. Zeigen Sie,
  dass $-A:=\{x\in\mathbb{R}\mid -x\in A\}$ nach oben beschränkt ist, und
  folgern Sie $\inf A=-\sup(-A)$.
}{%
  $s$ untere Schranke von $A$ $\Leftrightarrow$ $-s$ obere Schranke von
  $-A$. Nach 1.1.12 existiert $S:=\sup(-A)$, und $-S$ ist größte untere
  Schranke von $A$, d.\,h. $\inf A=-S=-\sup(-A)$.
}

% --- Abschnitt 1.2 Konvergenz --------------------------------------------
\lernabschnitt{1.2 Konvergenz}

\lernkarte[61211-1-sa-06]{%
  Übung 1.2.1\\[4pt]
  Seien $\varepsilon_1,\dots,\varepsilon_n>0$ und
  $\varepsilon:=\min\{\varepsilon_1,\dots,\varepsilon_n\}$. Zeigen Sie
  $\displaystyle U_\varepsilon(a)\subseteq\bigcap_{k=1}^{n}U_{\varepsilon_k}(a)$.
}{%
  Es gilt sogar Gleichheit: $\varepsilon\le\varepsilon_k$ für alle $k$
  $\Rightarrow$ $d(x,a)<\varepsilon\Rightarrow d(x,a)<\varepsilon_k$, also
  $U_\varepsilon(a)\subseteq U_{\varepsilon_k}(a)$; umgekehrt liegt $\varepsilon$
  selbst in $\{\varepsilon_k\}$.
}

\lernkarte[61211-1-sa-07]{%
  Übung 1.2.2\\[4pt]
  Seien $f=(a_k)$, $g=(b_k)$ konvergent. Beweisen Sie die Monotonie des
  Grenzwerts: $a_k\le b_k$ für fast alle $k$ $\Rightarrow$
  $\lim f\le\lim g$.
}{%
  Widerspruch: Annahme $a:=\lim f>b:=\lim g$, $m:=\tfrac{a+b}{2}$. Dann
  $a_k\in\,]m,\infty[$ und $b_k\in\,]-\infty,m[$ f.\,f.\,a.\ $k$, also
  $b_k<m<a_k$ -- Widerspruch zu $a_k\le b_k$.
}

\lernkarte[61211-1-sa-08]{%
  Übung 1.2.3\\[4pt]
  Sei $(a_k)$ reell. Zeigen Sie die Äquivalenz von\\
  (i) $\forall\varepsilon\,\exists n_0\,\forall k>n_0:\,|a_k-a_{n_0}|<\varepsilon$\\
  (ii) $\forall\varepsilon\,\exists n_0\,\forall k,\ell\ge n_0:\,|a_k-a_\ell|<\varepsilon$.
}{%
  \textbf{(i)$\Rightarrow$(ii):} mit $\varepsilon'=\tfrac{\varepsilon}{2}$
  ist $|a_k-a_\ell|\le|a_k-a_{n_0}|+|a_{n_0}-a_\ell|<\varepsilon$.\\
  \textbf{(ii)$\Rightarrow$(i):} setze $\ell=n_0$.
}

\lernkarte[61211-1-sa-09]{%
  Übung 1.2.4\\[4pt]
  Zeigen Sie, dass 1.2.5 auch für $a\in\{-\infty,\infty\}$ gilt: die
  Umgebungseigenschaften (i)--(iii) für eine Umgebung $U$ von $a$.
}{%
  (i) $a\in U$ nach Definition.\\
  (ii) $M\supseteq U\Rightarrow M$ Umgebung ($]\alpha,\infty[\subseteq U\subseteq M$).\\
  (iii) endlicher Durchschnitt: $\alpha:=\max\{\alpha_k\}$, dann
  $]\alpha,\infty[\subseteq\bigcap U_k$.
}

\lernkarte[61211-1-sa-10]{%
  Übung 1.2.5\\[4pt]
  Beweisen Sie die $\varepsilon$--$n_0$--Kriterien 1.2.25 für uneigentliche
  Grenzwerte $\lim a_k=\pm\infty$.
}{%
  \[ \lim_{k\to\infty}a_k=\infty \iff
     \forall\varepsilon>0\,\exists n_0\,\forall k\ge n_0:\,a_k>\tfrac{1}{\varepsilon}, \]
  analog $\lim a_k=-\infty \iff a_k<-\tfrac{1}{\varepsilon}$ (folgt aus den
  Definitionen der Umgebungen von $\pm\infty$).
}

% --- Abschnitt 1.3 R^n als reeller Vektorraum ----------------------------
\lernabschnitt{1.3 $\mathbb{R}^n$ als reeller Vektorraum}

\lernkarte[61211-1-sa-11]{%
  Übung 1.3.1\\[4pt]
  Beweisen Sie die Rechenregeln (Bemerkungen (5),(6) zu 1.3.3) in einem
  reellen Vektorraum $X$.
}{%
  Für $x,y,z\in X$, $\alpha\in\mathbb{R}$:
  \begin{itemize}[topsep=2pt,itemsep=1pt]
    \item $0x=0$ \quad(aus $0x+x=(0+1)x=x$)
    \item $(-1)x=-x$
    \item $x-(y+z)=(x-y)-z$
    \item $\alpha(x-y)=\alpha x-\alpha y$
  \end{itemize}
}

\lernkarte[61211-1-sa-12]{%
  Übung 1.3.2\\[4pt]
  Zeigen Sie, dass $\mathrm{Abb}(M,\mathbb{R})$, $C(M)$ und $\mathcal{P}_n$
  reelle Vektorräume sind.
}{%
  Vektorraum-Axiome punktweise nachrechnen; neutrales Element ist die
  Nullabbildung $\widehat{0}$. $C(M)$ und $\mathcal{P}_n$ sind Unterräume
  von $\mathrm{Abb}$ (abgeschlossen unter $+$ und $\alpha\cdot$).
}

\lernkarte[61211-1-sa-13]{%
  Übung 1.3.3\\[4pt]
  Seien $e_1,\dots,e_n$ die Standardbasisvektoren und
  $x={}^t(x_1,\dots,x_n)\in\mathbb{R}^n$. Zeigen Sie: (i) $x=\sum_{k=1}^n x_k e_k$;
  (ii) diese Darstellung ist eindeutig.
}{%
  (i) $\sum x_k e_k={}^t(x_1,\dots,x_n)=x$ (direkt komponentenweise).\\
  (ii) Aus $x=\sum\alpha_k e_k={}^t(\alpha_1,\dots,\alpha_n)$ folgt
  $\alpha_k=x_k$. Jedes $x$ ist also eindeutige Linearkombination der $e_k$.
}

\lernkarte[61211-1-sa-14]{%
  Übung 1.3.4\\[4pt]
  Seien $a,b\in\mathbb{R}^2$, $a\neq b$. Beschreiben Sie
  $G_{ab}:=\{a+t(b-a)\mid t\in\mathbb{R}\}$ geometrisch und als
  $\{\,{}^t(x,y)\mid Ax+By+C=0\}$.
}{%
  $G_{ab}$ ist die \textbf{Gerade durch $a$ und $b$}. Elimination von $t$
  liefert $Ax+By+C=0$ mit
  \[ A=-(b_2-a_2),\quad B=b_1-a_1,\quad C=a_1b_2-a_2b_1. \]
}

% --- Abschnitt 1.4 R^n als normierter Raum -------------------------------
\lernabschnitt{1.4 $\mathbb{R}^n$ als normierter Raum}

\lernkarte[61211-1-sa-15]{%
  Übung 1.4.1\\[4pt]
  Sei $(X,\|\ \|)$ normiert. (i) Beweisen Sie die zweite
  Dreiecksungleichung $\bigl|\,\|x\|-\|y\|\,\bigr|\le\|x+y\|$.
  (ii) Folgern Sie $|d(x,z)-d(z,y)|\le d(x,y)$.
}{%
  (i) $\|x\|=\|(x+y)-y\|\le\|x+y\|+\|y\|$ und symmetrisch
  $\Rightarrow\bigl|\,\|x\|-\|y\|\,\bigr|\le\|x+y\|$.\\
  (ii) $x'=x-z,\ y'=z-y$ in (i) einsetzen ($x'+y'=x-y$).
}

\lernkarte[61211-1-sa-16]{%
  Übung 1.4.2\\[4pt]
  Zeigen Sie, dass $d(x,y):=\dfrac{|x-y|}{1+|x-y|}$ ein Abstand auf
  $\mathbb{R}$ ist. Ist $d$ durch eine Norm erzeugt?
}{%
  Definitheit und Symmetrie klar; Dreiecksungleichung mit
  $h(t)=\tfrac{t}{1+t}$ monoton wachsend.\\
  \textbf{Nicht} normerzeugt: Homogenität verletzt
  (z.\,B.\ $\alpha=2,x=1$: $\tfrac{|\alpha x|}{1+|\alpha x|}\neq|\alpha|\tfrac{|x|}{1+|x|}$).
}

\lernkarte[61211-1-sa-17]{%
  Übung 1.4.3\\[4pt]
  Sei $\ast$ ein Skalarprodukt auf $X$ und $\|x\|:=\sqrt{x\ast x}$.
  Beweisen Sie (i) $(x\ast y)^2\le\|x\|^2\|y\|^2$; (ii) $\|\ \|$ ist eine Norm.
}{%
  (i) Aus $0\le(\alpha x+y)\ast(\alpha x+y)$ mit
  $\alpha=-\tfrac{x\ast y}{\|x\|^2}$ folgt Cauchy--Schwarz.\\
  (ii) Definitheit und Homogenität direkt; Dreiecksungleichung aus (i).
}

% --- Abschnitt 1.5 Konvergenz in R^n -------------------------------------
\lernabschnitt{1.5 Konvergenz in $\mathbb{R}^n$}

\lernkarte[61211-1-sa-18]{%
  Übung 1.5.1\\[4pt]
  Sei $x_k:={}^t\!\left(\dfrac{\ln k}{k},\dfrac{2^k}{k!},\sqrt[k]{k}\right)\in\mathbb{R}^3$.
  Zeigen Sie, dass $(x_k)$ in $\mathbb{R}^3$ konvergiert.
  (Hinweis: $\ln k<2\sqrt{k}$.)
}{%
  Koordinatenweise (1.5.11): $\tfrac{\ln k}{k}\to0$ (Sandwich, da
  $0\le\tfrac{\ln k}{k}<\tfrac{2}{\sqrt k}$), $\tfrac{2^k}{k!}\to0$
  (Reihenglied von $\exp2$), $\sqrt[k]{k}=e^{\ln k/k}\to1$. Grenzwert
  ${}^t(0,0,1)$.
}

\lernkarte[61211-1-sa-19]{%
  Übung 1.5.2\\[4pt]
  Sei $f_k:[0,1]\to\mathbb{R}$, $f_k(t):=t^k$. Konvergiert $(f_k)$
  (i) in $(C[0,1],\|\ \|_\infty)$, (ii) in $(C[0,1],\|\ \|_1)$?
}{%
  (i) \textbf{Nein} in $\|\ \|_\infty$: der punktweise Grenzwert
  ($0$ auf $[0,1[$, $1$ bei $t=1$) ist unstetig.\\
  (ii) \textbf{Ja} in $\|\ \|_1$ gegen $\widehat{0}$:
  $\|f_k\|_1=\int_0^1 t^k\,dt=\tfrac{1}{k+1}\to0$.
}

\lernkarte[61211-1-sa-20]{%
  Übung 1.5.3\\[4pt]
  Seien $\|\ \|$ und $\|\ \|^{*}$ zwei Normen auf $\mathbb{R}^n$. Zeigen Sie,
  dass es $\alpha,\beta>0$ gibt mit $\|x\|\le\alpha\|x\|^{*}$ und
  $\|x\|^{*}\le\beta\|x\|$.
}{%
  Alle Normen auf $\mathbb{R}^n$ sind äquivalent. Über $\|\ \|_\infty$ als
  Brücke (1.5.10): $\|x\|\le\alpha_1\|x\|_\infty$,
  $\|x\|_\infty\le\beta_2\|x\|^{*}$ usw.; setze
  $\alpha=\alpha_1\beta_2$, $\beta=\alpha_2\beta_1$.
}

\lernkarte[61211-1-sa-21]{%
  Übung 1.5.4\\[4pt]
  Beweisen Sie die Regeln 1.5.13: Aus $x_k\to a$ und $y_k\to b$ in
  $(X,\|\ \|)$ folgt $x_k+y_k\to a+b$ und $\alpha x_k\to\alpha a$.
}{%
  $\|(x_k+y_k)-(a+b)\|\le\|x_k-a\|+\|y_k-b\|\to0$ und
  $\|\alpha x_k-\alpha a\|=|\alpha|\,\|x_k-a\|\to0$.
}
